\section{EduSkillBench}
\label{sec:benchmark}
\label{sec:framework}

\subsection{Task Definition and Construction}
Given a teaching situation and request, an agent produces an instructional
artifact: for example, a lesson plan, diagnostic question, support sequence,
project brief, or practice schedule. Successful work requires interpreting
the available evidence, choosing appropriate teaching operations, and
delivering a correct, feasible product. Each task binds the situation and
request to a designated Skill and a fixed set of applicable rubric dimensions.
The Skill supplies reusable procedural support; it is not a reference answer.

The unified inventory contains 305 tasks across 14 Skills. It combines 263
Chinese teaching scenarios adapted from 30 documents in the archived
\texttt{edu\_scene\_66} collection with 42 constructed English tasks, three
per Skill. The source scenarios cover nine families; the constructed tasks
extend coverage to all 14. Versioned revisions repair identified mismatches
between questions, contexts, and criteria. Each constructed task specifies
the requested artifact and constraints at one of three nominal difficulty
levels. These levels organize construction, rather than constitute a
calibrated proficiency scale. All 305 tasks receive equal weight in the
primary result.

The packages are drawn from public educational Skill collections
\citep{educationagentskills,teachingskills}. They cover planning and alignment,
diagnostic questioning, adaptive support, learner participation, and practice
over time. Family sizes range from three to 157 tasks; lesson building
accounts for 51.48\% of the inventory. Figure~\ref{fig:skill-breakdown} reports
all family counts, and Appendix~\ref{app:inventory} preserves the full task
and training distribution.

\paragraph{An example of procedure fit.}
One task asks how to help a fourth-grade pupil label rock specimens independently
when labels are crooked or omit names and collection sites. Its adaptive-hint
Skill connects diagnosis to graduated assistance and an independence check.
The response must provide usable prompts and checks suited to this obstacle;
reproducing an entire teaching template is unnecessary. The applicable rubric
covers accuracy, scaffolding, guidance, age suitability, and safety. This
task--procedure--criterion binding makes the intended support observable.
It does not provide an answer key for every new example the agent may invent.

\subsection{Paired Agent Execution}
\label{sec:protocol}
The native protocol uses an isolated Docker workspace for each task and
OpenCode 1.18.11 with an educational single-turn agent instruction.
\noskill{} exposes no educational package. \withskill{} makes the designated
task-matched package available through the native Skill interface. The agent
decides whether to load it, inspect resources, and apply particular operations.
This measures autonomous use of the available candidate; it does not test
retrieval across all 14 Skills. Rubrics and evaluation references are withheld
from generation.

The agent deadline is 600 seconds, with a 480-second provider timeout and a
10,000-token output cap per model call. Multiple calls may occur within a
trajectory. Paired runs retain the requested artifact, available tool records,
and execution status. There is one accepted outcome per model--task--condition,
and completed answers are reused when only judging needs recovery.
Exact model settings and the retained GLM-5.3 exposure inventory are documented
in Appendix~\ref{app:execution}.

\subsection{Reference-Conditioned Artifact Evaluation}
\label{sec:judge}
DeepSeek V4.1 Flash jointly evaluates each candidate and the task's saved
DeepSeek V4 Pro \noskill{} answer. The two texts appear as anonymous A/B
responses, with hash-assigned order, the task, and its applicable rubric.
The evaluator sees the retained delivery text; files absent from that input
are not automatically inspected.
Following the comparison mechanism of QUEST \citep{quest}, both receive
dimension scores $s_{x,i,d}\in[0,10]$. We compute
\begin{align}
J_{x,i}&=\frac{1}{10D_i}\sum_{d=1}^{D_i}s_{x,i,d},\\
R_i&=\frac{J_{c,i}}{J_{c,i}+J_{r,i}},\qquad
S=\frac{100}{305}\sum_{i=1}^{305}R_i,
\label{eq:refscore}
\end{align}
where $D_i$ is the number of applicable dimensions and $c,r$ denote candidate
and reference. All tasks use this equal-weight rule, regardless of their
construction origin. $R_i=0.5$ means parity with the reference under that
judgment, not 50\% accuracy. Pro's reference-self condition receives 0.5;
execution failures receive zero in the fixed 305-task denominator. A pair
with both components zero is unresolved; none remain in the reported archive.

The judge checks correctness, consistency, feasibility, and task constraints;
scores of at least 9 require response evidence. Substantive errors reduce
relevant dimensions. Structural validators check coverage, evidence bindings,
and score ranges. The reference is a comparison anchor, not a gold answer:
generated scientific examples have no independent computational oracle.
We therefore retain both $J$ components and compare judgments with expert
adjudication in Section~\ref{sec:human-summary}. A stable ratio can conceal
shared changes in candidate and reference scores.
